\documentclass[10pt,a4paper]{amsart}
\usepackage[margin=19mm]{geometry}
\usepackage{amsmath,amssymb,mathtools}
\usepackage[scr=rsfs,cal=euler]{mathalfa}
\usepackage{xfrac}
\usepackage[unicode,hidelinks,bookmarks=false]{hyperref}
\newcommand{\ie}{i.e.\ }
\newcommand{\cf}{cf.\ }
\newcommand{\resp}{respectively\ }
\newcommand{\Subsection}{\subsection}
\providecommand{\nobreakdash}{\nobreak\hbox{-}}
\setlength{\emergencystretch}{2em}
\multlinegap0pt
\usepackage{bm}
\usepackage{bbm}
\usepackage{dsfont}
\usepackage{xspace}
\usepackage{cancel}
\usepackage{mathtools}
\usepackage[shortlabels]{enumitem}
\usepackage{amsthm}
\makeatletter
\@ifundefined{lemma}{\newtheorem{lemma}{Lemma}}{}
\makeatother
\newcommand{\R}{{\mathbb R}}
\newcommand{\betabf}{{\boldsymbol{\beta}}}
\newcommand{\fcnu}{{f^{d, s}_{c, \{\nu_{L, U}\}}}} 
\newcommand{\fstinf}{{\mathcal{F}^{d, s}_{\infty-\text{ST}}}}
\newcommand{\fstmid}{{\mathcal{F}^{d, s}_{\text{ST}, \bullet}}}  
\newcommand{\fst}{{\mathcal{F}^{d, s}_{\text{ST}}}}
\newcommand{\ind}{{\mathbf{1}}}
\newcommand{\onevec}{{\mathbf{1}}}
\newcommand{\pbf}{{\mathbf{p}}}
\newcommand{\qbf}{{\mathbf{q}}}
\newcommand{\vinfxgb}{{V^{d, s}_{\infty-\text{XGB}}}}
\newcommand{\xbf}{{\mathbf{x}}}
\newcommand{\zbf}{{\mathbf{z}}}
\title[What Functions Does XGBoost Learn?]{What Functions Does XGBoost Learn?}
\author{Dohyeong Ki, Adityanand Guntuboyina}
\date{Source: arXiv:2601.05444v1 (CC BY 4.0)}
\begin{document}
\maketitle

\noindent\emph{Source and licence.} What Functions Does XGBoost Learn?. Version arXiv:2601.05444v1 (CC BY 4.0), distributed under the Creative Commons Attribution 4.0 International licence, as recorded in the arXiv licence metadata for this exact version and shown on the abstract page https://arxiv.org/abs/2601.05444v1. Full rights notice: licenses/arxiv-2601.05444-CC-BY-4.0.txt.\par\smallskip

\noindent\textit{Editorial scope.} Counted unit: the Lemma whose source label is \texttt{lem:reduction-to-discrete-measures} in the article (source0.tex lines 3276--3290, source label \texttt{lem:reduction-to-discrete-measures}), delivered together with the article's own complete original proof of it (source0.tex lines 3293--3411, one \texttt{proof} environment of 119 lines). No mathematics is written, repaired or re-derived: statement and proof are the article's own text line for line. Required context retained uncounted: lem:discretization (lines 892--913); lattice (lines 898--903); eq:index-set (lines 3223--3228). Every definition, display and label of those lines is retained unchanged. Omitted from this package and not redistributed: 1--891, 904--3222, 3229--3275, 3291--3292, 3412--5766. Nothing inside a retained excerpt is omitted or altered: every retained body is delivered exactly as the article's lines are written, so no omission receipt of any kind accompanies this package. Citations: the retained excerpts carry no citation command at all, so no bibliography entry of the article is transcribed and no reference is redistributed. Reference adaptation: none was needed and none was made; every reference inside the delivered unit resolves inside it, so no unresolved reference or citation is produced. Printed slips kept literal: no printed word, symbol, hypothesis or number of the counted statement or of its proof was altered. Layout and class: the article's own class is \texttt{article} with packages including geometry, authblk, amsthm, amsmath, amssymb, natbib, hyperref, hypernat, mathabx, tikz, soul, multirow, float, enumitem; the wrapper is the lane's pinned \texttt{amsart} editorial wrapper at 10pt on A4, which restates the theorem-like environments the retained text needs and adds the article's own macro definitions that the retained text needs (\texttt{\textbackslash R}, \texttt{\textbackslash betabf}, \texttt{\textbackslash fcnu}, \texttt{\textbackslash fst}, \texttt{\textbackslash fstinf}, \texttt{\textbackslash fstmid}, \texttt{\textbackslash ind}, \texttt{\textbackslash onevec}, \texttt{\textbackslash pbf}, \texttt{\textbackslash qbf}, \texttt{\textbackslash vinfxgb}, \texttt{\textbackslash xbf}, \texttt{\textbackslash zbf}), with extra wrapper packages (bm, bbm, dsfont, xspace, cancel, mathtools, enumitem). The delivered numbering is the unit's own. Assets: no retained span contains a figure environment, a graphics-inclusion command, a PDF-page inclusion, a table or a TikZ picture; no image file is copied into this package, the package declares no asset dependency and no image file is redistributed with it. Licence: version arXiv:2601.05444v1 of this article is distributed under the Creative Commons Attribution 4.0 International licence, as recorded in the arXiv licence metadata for this exact version and shown on the abstract page https://arxiv.org/abs/2601.05444v1. A full rights notice is delivered at licenses/arxiv-2601.05444-CC-BY-4.0.txt. Characters: every retained excerpt is pure ASCII, so the delivered engine, XeLaTeX with its default Latin Modern text font, reports no missing character.
\begin{lemma}\label{lem:discretization}
    For every $f \in \fstinf$, there exists 
    $\fcnu \in \fstmid$ such that
    \begin{enumerate}[label = (\alph*)]
    \item $\nu_{L, U}$ are discrete signed Borel measures supported 
    on the lattices
    \begin{equation}\label{lattice}
        \prod_{j \in L} \Big\{ (v^{(j)}_1 + v^{(j)}_2)/2, \dots, 
        (v^{(j)}_{n_j - 1} + v^{(j)}_{n_j})/2 \Big\} \times 
        \prod_{j \in U} \Big\{ (v^{(j)}_1 + v^{(j)}_2)/2, \dots, 
        (v^{(j)}_{n_j - 1} + v^{(j)}_{n_j})/2 \Big\}
    \end{equation}
    \item $\fcnu(\xbf^{(i)}) = f(\xbf^{(i)})$ for 
    $i = 1, \dots, n$
    \item 
    \begin{equation*}
        \vinfxgb(\fcnu) 
        = \sum_{0 < |L| + |U| \le s} \|\nu_{L, U}\|_{\text{TV}}
        \le \vinfxgb(f).
    \end{equation*}
    \end{enumerate}
\end{lemma}

\begin{equation}\label{eq:index-set}
    J = \bigg\{ (L, U, \pbf, \qbf):  
    L, U \subseteq [d], 0 < |L| + |U| \le s,
    \pbf \in \prod_{j \in L} [n_j - 1], \text{ and } 
    \qbf \in \prod_{j \in U} [n_j - 1] \bigg\}
\end{equation}

\begin{lemma}\label{lem:reduction-to-discrete-measures}
    For every $\fcnu \in \fstinf$, there exists 
    $f^{d, s}_{b, \{\mu_{L, U}\}} \in \fst$ with discrete signed 
    measures $\mu_{L, U}$ supported on the lattices \eqref{lattice} 
    such that 
    \begin{enumerate}[label = (\alph*)]
    \item $f^{d, s}_{b, \{\mu_{L, U}\}}(\xbf^{(i)})
    = \fcnu(\xbf^{(i)})$ for $i = 1, \dots, n$
    \item 
    \begin{equation*}
        \sum_{0 < |L| + |U| \le s} \|\mu_{L, U}\|_{\text{TV}} 
        \le \sum_{0 < |L| + |U| \le s} \|\nu_{L, U}\|_{\text{TV}}.
    \end{equation*}
    \end{enumerate}
\end{lemma}

\begin{proof}[Proof of Lemma \ref{lem:discretization}]
For $z_1, \dots, z_n \in \R$, define 
\begin{equation*}
    \vinfxgb(z_1, \dots, z_n) = \inf \bigg\{ 
    \sum_{0 < |L| + |U| \le s} \|\nu_{L, U}\|_{\text{TV}}: 
    \fcnu(\xbf^{(i)}) 
    = z_i \text{ for } i = 1, \dots, n \bigg\}.
\end{equation*}
For simplicity, we suppress the dependence on the design points 
$\xbf^{(1)}, \dots, \xbf^{(n)}$.
By definition,
\begin{equation*}
    \vinfxgb(f(\xbf^{(1)}), \dots, f(\xbf^{(n)})) 
    \le \vinfxgb(f) \quad \text{for every } f \in \fstinf.
\end{equation*}

By Lemma \ref{lem:reduction-to-discrete-measures}, for each 
$z_1, \dots, z_n$, we have
\begin{equation*}
\begin{aligned}
    \vinfxgb(z_1, \dots, z_n) &= \inf \bigg\{ 
    \sum_{0 < |L| + |U| \le s} \|\nu_{L, U}\|_{\text{TV}}: 
    \fcnu(\xbf^{(i)}) 
    = z_i \text{ for } i = 1, \dots, n, \\
    &\qquad \qquad \qquad \qquad \qquad \qquad \qquad 
    \text{ with } \nu_{L, U} 
    \text{ supported on the lattices \eqref{lattice}}\bigg\}.
\end{aligned}
\end{equation*}
Recall the index set $J$ from \eqref{eq:index-set}, and let $\xbf$ 
denote the $n \times |J|$ matrix with entries
\begin{equation*}
    X_{i, (L, U, \pbf, \qbf)} 
    = \prod_{j \in L} 
    \ind\big(x^{(i)}_j \ge (v^{(j)}_{p_j} + v^{(j)}_{p_j + 1})/2\big)
    \cdot \prod_{j \in U} 
    \ind\big(x^{(i)}_j < (v^{(j)}_{q_j} + v^{(j)}_{q_j + 1})/2\big)
\end{equation*}
for $i = 1, \dots, n$ and $(L, U, \pbf, \qbf) \in J$.
We parametrize discrete signed measures $\nu_{L, U}$ 
supported on the lattices \eqref{lattice} by 
\begin{equation}\label{eq:measures-discrete-parametrization}
    \nu_{L, U} 
    \big(\{((v^{(j)}_{p_j} + v^{(j)}_{p_j + 1})/2, j \in L) 
    \times ((v^{(j)}_{q_j} + v^{(j)}_{q_j + 1})/2, j \in U)\}\big)
    = \beta^{L, U}_{\pbf, \qbf}
\end{equation}
for $L, U \subseteq [d]$ with $0 < |L| + |U| \le s$, 
$\pbf = (p_j, j \in L) \in \prod_{j \in L} [n_j - 1]$, 
and $\qbf = (q_j, j \in U) \in \prod_{j \in U} [n_j - 1]$.
With this parametrization, we can express $\vinfxgb(z_1, \dots, z_n)$ as 
\begin{equation}\label{eq:discrete-complexity}
    \vinfxgb(z_1, \dots, z_n) = \inf \big\{\|\betabf\|_1: \xbf \betabf 
    = \zbf - c \onevec \text{ for some } c \in \R\big\},
\end{equation}
where $\zbf = (z_1, \dots, z_n)$, and $\onevec$ is the all-ones 
vector.

Next, we show that if the set 
\begin{equation}\label{eq:discrete-feasible-set}
    \mathcal{D}_{\zbf} := \big\{\betabf \in \R^{|J|}: 
    \xbf \betabf = \zbf - c \onevec 
    \text{ for some } c \in \R\big\}
\end{equation}
is nonempty, which is clearly the case when
$\zbf = (f(\xbf^{(1)}), \dots, f(\xbf^{(n)}))$ 
for some $f \in \fstinf$, then there exists $\hat{\betabf}$ that achieves 
the minimum in \eqref{eq:discrete-complexity}.
Fix $\zbf \in \R^n$ such that $\mathcal{D}_{\zbf}$ is 
nonempty, and suppose $\xbf \betabf_0 = \zbf - c_0 \onevec$ 
for some $\betabf_0$ and $c_0$.
Clearly, the infimum on the right-hand side of 
\eqref{eq:discrete-complexity} remains unchanged if we further 
constrain $\betabf$ to satisfy $\|\betabf\|_1 \le \|\betabf_0\|_1$:
\begin{equation*}
    \vinfxgb(z_1, \dots, z_n) = \inf \big\{\|\betabf\|_1: \xbf \betabf 
    = \zbf - c \onevec \text{ for some } c \in \R 
    \text{ and } \|\betabf\|_1 \le \|\betabf_0\|_1\big\}.
\end{equation*}
It is also straightforward to verify that the set
\begin{equation*}
    \mathcal{D}_{\zbf} \cap \{\betabf \in \R^{|J|}: 
    \|\betabf\|_1 \le \|\betabf_0\|_1\}
\end{equation*}
is nonempty, closed, and bounded.
Since the map $\betabf \mapsto \|\betabf\|_1$ is continuous, it follows 
that there exists $\hat{\betabf}$ that attains the minimum in
\eqref{eq:discrete-complexity}.

Using the results established above, we now prove the lemma. 
Fix $f \in \fstinf$, and let $\hat{\betabf}$ be the minimizer of 
\eqref{eq:discrete-complexity} for $\zbf = (f(\xbf^{(1)}), \dots, 
f(\xbf^{(n)}))$.
Let $\hat{c}$ be the corresponding constant from 
\eqref{eq:discrete-feasible-set}, and let 
$\nu_{L, U}$ denote the signed Borel measures associated with 
$\hat{\betabf}$ via \eqref{eq:measures-discrete-parametrization}.
By construction, $\fcnu \in \fst$ 
satisfies the first two conditions of the lemma. 
Moreover, since
\begin{equation*}
\begin{aligned}
    &\vinfxgb(f(\xbf^{(1)}), \dots, f(\xbf^{(n)})) 
    = \|\hat{\betabf}\|_1 
    = \sum_{0 < |L| + |U| \le s} \|\nu_{L, U}\|_{\text{TV}} \\
    &\qquad \ge \vinfxgb(\fcnu)
    \ge \vinfxgb(f(\xbf^{(1)}), \dots, f(\xbf^{(n)})),
\end{aligned}
\end{equation*} 
we have
\begin{equation*}
    \vinfxgb(\fcnu) 
    = \sum_{0 < |L| + |U| \le s} \|\nu_{L, U}\|_{\text{TV}}
    = \vinfxgb(f(\xbf^{(1)}), \dots, f(\xbf^{(n)})) 
    \le \vinfxgb(f).
\end{equation*}
Hence, $\fcnu \in \fstmid$ is a function that satisfies 
all the desired properties.
\end{proof}

\end{document}
